\begin{proposition}[7.3.16]
\label{IV.2.s7.3.16-fr}
Se supone que la propiedad $\boldsymbol P$ es del primer \textup{(resp.} del
segundo) tipo, definida a partir de una propiedad $\boldsymbol R$
\textup{(7.3.10)}. Cuando $\boldsymbol P$ es del segundo tipo, se supone además
que la condición
\begin{quote}
«para toda extensión \textup{finita} $k'$ de $k$ y todo
$z'\in Z\otimes_k k'$, se cumple
$\boldsymbol R(\mathscr O_{z'})$»
\end{quote}
implica $\boldsymbol P(Z,k)$ \textup{(lo cual se verifica en los ejemplos
(iv$'$) a (vii$'$) de (7.3.8), en virtud de (6.7.7) y (4.6.1))}.

Sea $A$ un anillo semilocal noetheriano. Si la propiedad
$\boldsymbol R$ satisface las condiciones $(\mathrm R_{\mathrm I})$
y $(\mathrm R_{\mathrm{II}})$, las propiedades siguientes son
equivalentes:
\begin{enumerate}
\item[\emph{a)}] $A$ es un $\boldsymbol P$-anillo.
\item[\emph{b)}] Para todo anillo cociente íntegro $B$ de $A$
\textup{(resp.} toda $A$-álgebra finita e íntegra $B$) y todo ideal primo
$\mathfrak q$ de $\widehat B$ cuya imagen inversa en $B$ es 0, se
cumple $\boldsymbol R(\widehat B_{\mathfrak q})$.
\end{enumerate}
Si además $\boldsymbol R$ satisface la condición
$(\mathrm R'_{\mathrm I})$, las propiedades \textup{a)} y \textup{b)} son también
equivalentes a:
\begin{enumerate}
\item[\emph{c)}] Para todo anillo cociente $B$ de $A$ \textup{(resp.} toda
$A$-álgebra finita $B$), si se pone
$Y=\operatorname{Spec}(B)$,
$Y'=\operatorname{Spec}(\widehat B)$ y $g:Y'\to Y$ es el morfismo
canónico, se tiene (con las notaciones de \textup{(7.3.12)})
\begin{equation}
U_{\mathbf{R}}(Y')=g^{-1}(U_{\mathbf{R}}(Y)).
\tag{7.3.16.1}
\label{IV.2.eq.7.3.16.1-fr}
\end{equation}
\end{enumerate}
\end{proposition}

\begin{proof}
La equivalencia entre \emph{a)} y \emph{b)} es inmediata cuando $\boldsymbol P$ es
del primer tipo: en efecto, para todo ideal primo $\mathfrak p$ de
$A$, se ha visto que la fibra formal de $B=A/\mathfrak p$ en el punto
genérico (0) de $\operatorname{Spec}(B)$ no es sino la fibra formal de
$A$ en el punto $\mathfrak p\in\operatorname{Spec}(A)$ (7.3.13).

Cuando $\boldsymbol P$ es del segundo tipo, la equivalencia entre \emph{a)}
y \emph{b)} se deduce del siguiente lema más general:
\end{proof}

\begin{lemma}[7.3.16.2]
\label{IV.2.s7.3.16.2-fr}
Sean $A$ y $A'$ dos anillos, $\varphi:A\to A'$ un homomorfismo y
$f:\operatorname{Spec}(A')\to\operatorname{Spec}(A)$ el morfismo
correspondiente. Para que, para todo $x\in\operatorname{Spec}(A)$,
toda

\oldpage[IV-2]{197}

extensión finita $k$ de $\boldsymbol{k}(x)$ y todo punto $z$ de
$f^{-1}(x)\otimes_{\mathbf{k}(x)}k=Z$, se cumpla
$\boldsymbol R(\mathscr O_{Z,z})$, es necesario y suficiente que se
satisfaga la siguiente condición: para toda $A$-álgebra íntegra
finita $B$, si
$g:\operatorname{Spec}(A'\otimes_A B)\to\operatorname{Spec}(B)$ es
el morfismo deducido de $f$ por extensión a $B$ del anillo de base,
y si $T=g^{-1}(\xi)$ es la fibra del punto genérico $\xi$ de
$\operatorname{Spec}(B)$, entonces, para todo $t\in T$, se cumple
$\boldsymbol R(\mathscr O_{T,t})$.
\end{lemma}

\begin{proof}
La condición es trivialmente necesaria porque, si $x$ es el punto de
$\operatorname{Spec}(A)$ sobre el que se encuentra $\xi$, $\boldsymbol{k}(\xi)$
es una extensión finita de $\boldsymbol{k}(x)$. Recíprocamente, consideremos un
punto arbitrario $x\in\operatorname{Spec}(A)$ y sea $k$ una extensión
finita de $\boldsymbol{k}(x)$. Si se pone $\mathfrak p=\mathfrak j_x$, $\boldsymbol{k}(x)$ es
el cuerpo de fracciones de $A/\mathfrak p$, y existe una base de $k$
sobre $\boldsymbol{k}(x)$ formada por elementos enteros sobre $A/\mathfrak p$;
si $B$ es el subanillo de $k$ engendrado por estos elementos, $B$ es
una $A$-álgebra íntegra finita y $k$ es el cuerpo $\boldsymbol{k}(\xi)$ en el
punto genérico $\xi$ de $\operatorname{Spec}(B)$; como $x$ es la
imagen de $\xi$ en $\operatorname{Spec}(A)$, la fibra
$g^{-1}(\xi)$ no es sino $f^{-1}(x)\otimes_{\mathbf{k}(x)}k$, lo que termina
la demostración del lema.
\end{proof}

\begin{proof}[Fin de la demostración de (7.3.16)]
Que \emph{a)} implique \emph{c)} es una consecuencia inmediata de (7.3.15) y
(7.3.12). Por otra parte, si se especializa \emph{c)} al caso en que $B$
es íntegro, y $y$ designa el punto genérico de
$Y=\operatorname{Spec}(B)$, se tiene $\mathscr O_{Y,y}=\boldsymbol{k}(y)$; por
tanto, $y\in U_{\mathbf{R}}(Y)$, puesto que $\boldsymbol R(k)$
se cumple para todo cuerpo $k$. Al expresar que todo punto
$y'\in g^{-1}(y)$ pertenece a $U_{\mathbf{R}}(Y')$, se obtiene el enunciado \emph{b)}; así, \emph{c)} implica \emph{b)}. Queda demostrado.
\end{proof}
